\section{怎样发现好 Filter：不要相信直觉，先跑 Ablation}

“清洗数据”最容易被写成规则清单，但真正困难的是判断规则是否提升模型。课程把 filter discovery 变成实验科学：manual inspection 提出候选，small-model ablation 估计因果效果，多个 seeds 和 high-signal benchmarks 降低噪声。

\subsection{从质量主张到可执行 pipeline}

第一张 evidence slide 引用 Yi/Open Foundation Models 的观点：standard architecture 加高质量数据也能表现很强。它只提供动机；下一页才展示 language filtering、quality rules、dedup、semantic/topic 等可执行 stages。

\lecturefigure{slide-39.jpg}{为什么要过滤：高质量 data 对标准 architecture 的影响}{Loubna Ben Allal 官方 deck，第 39 页}

\lecturefigure{slide-40.jpg}{通用 pretraining data-cleaning pipeline}{Loubna Ben Allal 官方 deck，第 40 页}

常见 rule/perplexity filter 可写成保留条件
\[
\mathbb{1}\{q(x)\in[a,b]\}\cdot \mathbb{1}\{\operatorname{PPL}_{r}(x)<\tau\},
\]
其中 $q(x)$ 是长度、重复行、符号比例等 heuristic，$r$ 是 reference model。阈值 $a,b,\tau$ 必须用样本和 ablation 校准。

\begin{warningbox}{Perplexity filter 的双刃剑}
高 PPL 可能是垃圾，也可能是稀有语言、代码、公式或新知识；低 PPL 可能是模板与重复。使用 reference-model PPL 会把其 bias 带入 dataset，应与 diversity/coverage 指标联合。
\end{warningbox}

\subsection{寻找最佳过滤技术：manual inspection 只是起点}

前一节列出 language、rule、perplexity 与 dedup stages，但还没有回答阈值怎样选。本节把 filter 设计变成可证伪实验：先看数据并解释异常，再训练 matched-compute 小模型验证；评测选择 early-training 就有 signal 的 benchmark，并用不同 seeds 避免把 optimizer noise 当成 filter gain。

\lecturefigure{slide-42.jpg}{Filter discovery：人工检查、小模型 ablation、高 signal benchmark 与多 seeds}{Loubna Ben Allal 官方 deck，第 42 页}

若 filter $f$ 与 baseline 的分数差为 $\Delta_s$，跨 $k$ 个 seeds 的均值和标准误可写为
\[
\bar\Delta=\frac{1}{k}\sum_{s=1}^{k}\Delta_s,
\qquad
\operatorname{SE}(\bar\Delta)=\frac{\operatorname{sd}(\Delta_s)}{\sqrt{k}}.
\]
只有 effect 超过 noise，才值得把规则应用到 TB-scale corpus。

\begin{lstlisting}[caption={数据过滤 ablation 的最小实验协议}]
for filter_candidate in candidates:
    subset = apply_filter(raw_subset, filter_candidate)
    for seed in seeds:
        model = train_small_model(subset, matched_tokens, seed)
        scores[filter_candidate, seed] = evaluate(high_signal_suite)
compare_mean_variance_cost_and_retained_diversity(scores)
\end{lstlisting}

\teachervoice{00:27:10--00:30:00，讲者给出两个负结果：按 comment density 过滤几乎没提升；删除少于 5 stars 的仓库移除 70\% 以上数据，并训练出最差 ablation。直觉必须让位于实验。}

\subsection{FineWeb 的 200+ ablations：数据 recipe 也需要 search budget}

FineWeb 团队训练 200 多个约 1B 参数、约 30B tokens 的 ablation models 来选择 filters。这个规模说明 data research 不是廉价预处理；它需要 model-training budget、benchmark harness 和统计纪律。

\lecturefigure{slide-43.jpg}{FineWeb：用 200+ ablation models 搜索 filtering recipe}{Loubna Ben Allal 官方 deck，第 43 页}

\begin{importantbox}{把 data recipe 当作 hyperparameter search}
Filter type、threshold、dedup radius、language mixture、source weights 与 sampling temperature 都是超参数。与 architecture search 一样，应避免 test-set tuning，保留 held-out evaluation，并记录失败实验。
\end{importantbox}

\subsection{本章小结}

好 filter 不是规则看起来合理，而是在 matched compute、多 seeds、多个 early-signal benchmarks 下提高模型，同时保留 diversity 与 governance 目标。Star filtering 的失败是本讲最重要的实践经验之一。

